\documentclass[tikz,border=4pt]{standalone}
\usepackage{ctex}
\usepackage{amsmath}
\usetikzlibrary{calc, arrows.meta, backgrounds}

\begin{document}
\begin{tikzpicture}[
  x={(1.8cm,0cm)},
  y={(0cm,1.8cm)},
  z={(-0.52cm,-0.42cm)},
  line cap=round,
  vec/.style={-{Stealth[length=7pt, width=4pt]}, thick},
  vecred/.style={-{Stealth[length=7pt, width=4pt]}, thick, red},
  vecblue/.style={-{Stealth[length=7pt, width=4pt]}, thick, blue},
]

% 共面向量示意：p = x*a + y*b（三向量共面）
% 共面平面用平行四边形轮廓示意

% 原点 O
\coordinate (O) at (0,0,0);

% 基向量 a（沿平面方向1），b（沿平面方向2，带轻微 z 偏移令其看起来在同一平面）
\coordinate (A) at (1.2,0,0);     % 方向 a（纯 x）
\coordinate (B) at (0.4,0.9,0);   % 方向 b（xy 平面内，非平行 a）

% 分解系数 x=0.8, y=0.7 → P = 0.8*A + 0.7*B
\coordinate (xA) at (0.96,0,0);         % x*a = 0.8*a
\coordinate (yB) at (0.28,0.63,0);      % y*b = 0.7*b
\coordinate (P)  at (1.24,0.63,0);      % p = x*a + y*b

% 平面网格辅助线（虚线，灰色）——展示共面
% 过 xA 画平行 b 的线，过 yB 画平行 a 的线
\coordinate (xA_yB) at ($(xA)+(yB)$);   % 就是 P
\draw[gray, dashed] (xA) -- (xA_yB);
\draw[gray, dashed] (yB) -- (xA_yB);
\draw[gray, dashed] (O)  -- (xA);
\draw[gray, dashed] (O)  -- (yB);

% 平面轮廓（浅灰色填充）
\begin{scope}[on background layer]
  \fill[gray!15] (O) -- (A) -- ($(A)+(B)$) -- (B) -- cycle;
\end{scope}
\draw[gray] (O) -- (A) -- ($(A)+(B)$) -- (B) -- cycle;

% 基向量 a, b
\draw[vec, black] (O) -- (A);
\node[below right] at ($(O)!0.6!(A)$) {$\vec{a}$};

\draw[vecblue] (O) -- (B);
\node[left] at ($(O)!0.6!(B)$) {$\vec{b}$};

% 向量 p（红色）
\draw[vecred] (O) -- (P);
\node[above right] at ($(O)!0.65!(P)$) {$\vec{p}$};

% 分量标注
\node[below, font=\small] at ($(O)!0.5!(xA)$) {$x\vec{a}$};
\node[right, font=\small] at ($(xA)!0.5!(P)$) {$y\vec{b}$};

% 原点标注
\node[below left] at (O) {$O$};

% 说明文字
\node[below] at (0.62,-0.22,0) {$\vec{p} = x\vec{a}+y\vec{b}$\quad（$\vec{a},\vec{b}$ 不共线，三向量共面）};

\end{tikzpicture}
\end{document}
